\documentclass[11pt]{article}
\usepackage{progtable}

\begin{document}
\progheader{Troisième}{3e}

\begin{proglist}

\seq{3C01}{Identités remarquables}{ALG}{Maîtriser les identités remarquables}{
\item 3C01-ALG01 — Développer (a+b)² et (a−b)²
\item 3C01-ALG02 — Factoriser avec les identités remarquables
\item 3C01-ALG03 — Simplifier des expressions complexes
}

\seq{3C02}{Systèmes d’équations linéaires}{ALG}{Résoudre des systèmes 2×2}{
\item 3C02-ALG04 — Résoudre par substitution
\item 3C02-ALG05 — Résoudre par combinaison linéaire
\item 3C02-ALG06 — Modéliser un problème par un système
}

\seq{3C03}{Fonctions — notion et représentation}{FUNC}{Introduire la notion de fonction}{
\item 3C03-FUNC01 — Comprendre la notion de fonction
\item 3C03-FUNC02 — Lire et tracer une courbe représentative
\item 3C03-FUNC03 — Distinguer fonction linéaire et affine
}

\seq{3C04}{Géométrie — Thalès et similitude}{GEO}{Généraliser les théorèmes en espace}{
\item 3C04-GEO01 — Appliquer le théorème de Thalès
\item 3C04-GEO02 — Utiliser la similitude dans les triangles
\item 3C04-GEO03 — Appliquer Pythagore en 3D
}

\seq{3C05}{Probabilités — calculs et arbres}{STAT}{Raisonner sur des événements aléatoires}{
\item 3C05-STAT01 — Calculer une probabilité par dénombrement
\item 3C05-STAT02 — Construire et lire un arbre de probabilités
\item 3C05-STAT03 — Aborder les probabilités conditionnelles simples
}

\seq{3C06}{Révisions brevet — résolution de problèmes}{BREV}{Préparer le DNB en situation complexe}{
\item 3C06-BREV01 — Mobiliser plusieurs notions dans un problème
\item 3C06-BREV02 — Rédiger une démonstration complète
\item 3C06-BREV03 — Communiquer et argumenter mathématiquement
}

\end{proglist}
\end{document}
